% Quelle des Templates: https://de.wikipedia.org/wiki/LaTeX

\documentclass[12pt]{exam}
\usepackage[T1]{fontenc}
\usepackage{lmodern}
\usepackage[ngerman]{babel}
\usepackage{amsmath}
\usepackage{graphicx}
\usepackage{tabularx}
\usepackage{tikz}
\usepackage{pgfplots}
\usepackage{amssymb}
\usepackage{trfsigns}
\usetikzlibrary{positioning}
\usepackage{hyperref}
\usepackage{verbatim}
\usepackage{listings}
\usepackage[siunitx,european]{circuitikz}
\usepackage{subcaption}
\usepackage{pythonhighlight}
\usepackage[utf8]{inputenc}

\newcommand{\emptybox}[2][\textwidth]{%
  \begingroup
  \setlength{\fboxsep}{-\fboxrule}%
  \noindent\framebox[#1]{\rule{0pt}{#2}}%
  \endgroup
}

\begin{document}

\begin {center}
Technische Hochschule\\
Würzburg-Schweinfurt \\
Faculty Electrical Engineering\\
Summer 2026 \\
\end{center}

\begin{tabular}{ll}
Sprachsteuerung:& Degree program Elektro- und Informationstechnik (Bachelor) \\
Prozessdatenverarbeitung:& Degree program Elektro- und Informationstechnik (Bachelor) \\
Speech Recognition:& Degree program Robotics (Bachelor) \\
&\\
Date: & 07.05.2026 \\
Duration: & 90 Minutes \\
Auxiliary means: & one page DIN A4 formulary (both sided) and calculator \\
Number of tasks: & 5\\
&\\
\end{tabular}

\begin{tabular}{ll}
Name:&...........................................\\
&\\
Matriculation number:&...........................................\\
&\\
Degree program:&...........................................\\
&\\
First examiner:&Prof. Dr.-Ing. Martin Spiertz\\
Second examiner:&\\
&\\
\end{tabular}

\begin{tabular}{l}
Important hints: \\
Complete the tasks exclusively in the respective spaces provided in the worksheet.\\
If additional space is required, supplementary sheets will be distributed.\\
Please label all supplementary sheets with your name.
\end {tabular}

\begin{questions}
\clearpage %Aufgabe 1
\question Task (10 points):\\
\begin{parts}
\part Give the mathematical description of a digital signal, which has the highest possible digital level in dB FS. Evaluate the corresponding level in dB FS.\\\\
\begin{centering}
\begin{tikzpicture} %Bild 3
\draw[step =.5cm ,gray ,very  thin] (0,0)  grid(15 ,5);
\end{tikzpicture}
\end{centering}\\\\
A signal $x(t)=\cos\left(500\cdot t\right)$ is sampled with $r=249$ Hz.
\part Is this sampling causing Alias?\\\\
\begin{centering}
\begin{tikzpicture} %Bild 3
\draw[step =.5cm ,gray ,very  thin] (0,0)  grid(15 ,2);
\end{tikzpicture}
\end{centering}\\\\
\part Evaluate the sampled signal for $n=2$. Hint: $x(n)=x\left(t=n/r\right)$\\\\
\begin{centering}
\begin{tikzpicture} %Bild 3
\draw[step =.5cm ,gray ,very  thin] (0,0)  grid(15 ,7);
\end{tikzpicture}
\end{centering}\\\\
\clearpage
A FIR filter is given by: $h(n)=\left[1; -1\right]$
\part Is $h(n)$ stable? (give a reason for your answer)\\\\
\begin{centering}
\begin{tikzpicture} %Bild 3
\draw[step =.5cm ,gray ,very  thin] (0,0)  grid(15 ,2);
\end{tikzpicture}
\end{centering}\\\\
\part Evaluate the real part of the transfer function $H(f)$ of $h(n)$ with the z-Transform. Is the real part negative for any $f$?\\\\
\begin{centering}
\begin{tikzpicture} %Bild 3
\draw[step =.5cm ,gray ,very  thin] (0,0)  grid(15 ,5);
\end{tikzpicture}
\end{centering}\\\\
A signal $y(n)$ is generated through the following Python code:
\begin{python}
import numpy as np
y = np.random.randn(10000)
\end{python}
\part Give the Python code to evaluate the mean value and the standard deviation of a block of $y(n)$ starting at $n=1000$ with a blocklength of $N=1000$.\\\\
\begin{centering}
\begin{tikzpicture} %Bild 3
\draw[step =.5cm ,gray ,very  thin] (0,0)  grid(15 ,5);
\end{tikzpicture}
\end{centering}\\\\
\end{parts}

\clearpage %Aufgabe 2
\question Task (10 points):\\
\begin{parts}
\part Give two reasons for using zeropadding for a DFT.\\\\
\begin{centering}
\begin{tikzpicture} %Bild 3
\draw[step =.5cm ,gray ,very  thin] (0,0)  grid(15 ,5);
\end{tikzpicture}
\end{centering}\\\\
\part Why should a window function be multiplied with a signal before applying a DFT?\\\\
\begin{centering}
\begin{tikzpicture} %Bild 3
\draw[step =.5cm ,gray ,very  thin] (0,0)  grid(15 ,5);
\end{tikzpicture}
\end{centering}\\\\
\clearpage
In the following a STFT is applied to the signal $x(n)$. The sampling rate is $r=48$ kHz, the signal has a duration of $70$ minutes. The blocklength is $N=1000$, the overlapping is $75$ \%, the transform length is $K=2028$.
\part What are the advantages of using a large blocklength $N$? What are the advantages of using a small blocklength $N$?\\\\
\begin{centering}
\begin{tikzpicture} %Bild 3
\draw[step =.5cm ,gray ,very  thin] (0,0)  grid(15 ,5);
\end{tikzpicture}
\end{centering}\\\\
\part Determine the size of the matrix of the resulting spectrogram. \\\\
\begin{centering}
\begin{tikzpicture} %Bild 3
\draw[step =.5cm ,gray ,very  thin] (0,0)  grid(15 ,7);
\end{tikzpicture}
\end{centering}\\\\
\part What is the meaning of the abbreviations DFT, FFT and STFT?
\\\\
\begin{centering}
\begin{tikzpicture} %Bild 3
\draw[step =.5cm ,gray ,very  thin] (0,0)  grid(15 ,2);
\end{tikzpicture}
\end{centering}
\end{parts}

\clearpage %Aufgabe 3
\question Task (10 points):\\
\begin{parts}
\part A signal has the frequency $f=1234$ Hz. Determine the corresponding frequencies in Bark and mel.
\\\\
\begin{centering}
\begin{tikzpicture} %Bild 3
\draw[step =.5cm ,gray ,very  thin] (0,0)  grid(15 ,5);
\end{tikzpicture}
\end{centering}\\\\
\part What is the time resolution in milliseconds and the frequency resolution in Bark for the human ear?
\\\\
\begin{centering}
\begin{tikzpicture} %Bild 3
\draw[step =.5cm ,gray ,very  thin] (0,0)  grid(15 ,5);
\end{tikzpicture}
\end{centering}\\\\
\clearpage
In the following, it is assumed that $y(n)=x(n)+d(n)$, $x(n)$ und $d(n)$ are uncorrelated. $d(n)$ is additive white Gaussian noise.
\part For which $p$ the following is true: $\left|Y(f)\right|^p\approx\left|X(f)\right|^p+\left|D(f)\right|^p$\\\\
\begin{centering}
\begin{tikzpicture} %Bild 3
\draw[step =.5cm ,gray ,very  thin] (0,0)  grid(15 ,2);
\end{tikzpicture}
\end{centering}\\\\
\part What is the abbreviation of this assumed noise model?
\\\\
\begin{centering}
\begin{tikzpicture} %Bild 3
\draw[step =.5cm ,gray ,very  thin] (0,0)  grid(15 ,3);
\end{tikzpicture}
\end{centering}\\\\
\part Which property is fullfilled by the spectrum of  $d(n)$?\\\\
\begin{centering}
\begin{tikzpicture} %Bild 3
\draw[step =.5cm ,gray ,very  thin] (0,0)  grid(15 ,3);
\end{tikzpicture}
\end{centering}\\\\
\part Which properties are fullfilled by the samples of  $d(n)$ in time domain?\\\\
\begin{centering}
\begin{tikzpicture} %Bild 3
\draw[step =.5cm ,gray ,very  thin] (0,0)  grid(15 ,5);
\end{tikzpicture}
\end{centering}
\end{parts}

\clearpage %Aufgabe 4
\question Task (10 points):\\
\begin{parts}
\part Human speech is recorded by a microphone. What is the typical frequency range of human speech? What is the typical mean value of the recorded signal?
\\\\
\begin{centering}
\begin{tikzpicture} %Bild 3
\draw[step =.5cm ,gray ,very  thin] (0,0)  grid(15 ,6);
\end{tikzpicture}
\end{centering}\\\\
\part Which delay in milliseconds is produced by a symmetric FIR-filter in the passband? The length of the filter is $N=20$, the sampling rate is $r=16$ kHz. 
\\\\
\begin{centering}
\begin{tikzpicture} %Bild 3
\draw[step =.5cm ,gray ,very  thin] (0,0)  grid(15 ,6);
\end{tikzpicture}
\end{centering}\\\\
\clearpage
In the following, the IIR-filter  $y(n)=a\cdot y(n-1)+(1-a)\cdot x(n)$ is discussed.
\part For which value of $a$ is the filter stable?
\\\\
\begin{centering}
\begin{tikzpicture} %Bild 3
\draw[step =.5cm ,gray ,very  thin] (0,0)  grid(15 ,3);
\end{tikzpicture}
\end{centering}\\\\
\part Give the gain of the DC component:  $H\left(f=0\right)$.\\\\
\begin{centering}
\begin{tikzpicture} %Bild 3
\draw[step =.5cm ,gray ,very  thin] (0,0)  grid(15 ,5);
\end{tikzpicture}
\end{centering}\\\\
\part An input signal $x(n)$ is given. Write Python code to implement the above given IIR-filter.\\\\
\begin{python}
import numpy as np
a = 0.9
N = 1000
x = np.random.randn(N)
\end{python}
\begin{centering}
\begin{tikzpicture} %Bild 3
\draw[step =.5cm ,gray ,very  thin] (0,0)  grid(15 ,5);
\end{tikzpicture}
\end{centering}
\end{parts}

\clearpage %Aufgabe 5
\question Task (10 points):\\
For all ten following Python-snippets the following introduction line is assumed:
\begin{python}
import numpy as np
\end{python}
Determine for each snippet  if the ouput is A or B (by marking the correct output on the sheet).

\begin{python}
x = np.ones(5)
assert x.sum()==5, `A'
print(`B')
\end{python}

\begin{python}
x = np.array([1, 0, 0])
a = np.array([1, -0.5])
y = np.array([0, 0])
for i in range(2):
    y[i] = x[i] + 0.5 * y[i-1] if i > 0 else x[i]
assert (y == [1, 0.5]).all(), `A'
print(`B')
\end{python}

\begin{python}
x = np.array([1, 0, 0])
b = np.array([1, 1, 1])
y = np.convolve(x, b, mode='full')
assert (y == [1, 1, 1, 0, 0]).all(), `A'
print(`B')
\end{python}

\begin{python}
x = np.array([1, 0, 0, 0])
X = np.fft.fft(x)
assert X[0] == X[-1], `A'
print(`B')
\end{python}

\begin{python}
x = np.array([1, 2, 3])
b = np.array([1, 1])
y = np.convolve(x, b, mode='valid')
assert (y == [3, 5]).all(), `A'
print(`B')
\end{python}

\begin{python}
x = np.array([1, 2, 3, 4])
X = np.fft.fft(x)
assert np.isreal(X).all(), `A'
print(`B')
\end{python}

\begin{python}
x = np.array([1, 1, 1, 1])
X = np.fft.fft(x)
assert abs(X[1]) == 0, `A'
print(`B')
\end{python}

\begin{python}
x = np.array([1, 2, 3])
assert np.var(x) == 0.6666666666666666, `A'
print(`B')
\end{python}

\begin{python}
x = np.array([1, 2, 3])
assert np.std(x) == 1, `A'
print(`B')
\end{python}

\begin{python}
x = np.array([1, 2, 3])
assert np.mean(x) == 2, `A'
print(`B')
\end{python}

\end{questions}
\end{document}



